\section{ماژول kronecker\_compressor (Kronecker Product Compressor)}
\label{sec:module_kronecker_compressor}

\subsection{معرفی و انگیزه}
این مکانیزم بر مبنای مبانی نظری پیشرفته در ادبیات توجه و فشرده‌سازی \cite{eckart1936approximation} توسعه یافته است.

کاهش پارامترهای لایه‌های خطی در ترنسفورمرها معمولاً به تجزیه‌های کم‌رتبه ماتریسی همچون \lr{SVD} محدود می‌شود؛ اما این تجزیه‌ها در رتبه‌های بسیار پایین با افت چشمگیر ظرفیت بازنمایی مواجه می‌شوند، زیرا یک ماتریس با ابعاد $M \times N$ را به دو ماتریس $M \times r$ و $r \times N$ تبدیل می‌کنند که تنها زیرفضای خطی ساده‌ای را تقریب می‌زند.

ماژول \lr{kronecker\_compressor} از ضرب تانسوری کرونکر (\lr{Kronecker Product}) برای تجزیه ماتریس‌های وزن بهره می‌برد. ضرب کرونکر اجازه می‌دهد تا یک ماتریس بزرگ از وزن‌ها به صورت حاصل‌ضرب کرونکر دو ماتریس فاکتور بسیار کوچک‌تر نمایش داده شود. این ساختار نه تنها کاهش رادیکالی در تعداد پارامترها و حافظه ذخیره‌سازی به همراه دارد، بلکه به دلیل ویژگی‌های جبری حاصل‌ضرب تانسوری، انجام ضرب ماتریس-بردار را بدون تشکیل صریح ماتریس بزرگ بازسازی‌شده (\lr{Memory Materialization}) و با هزینه محاسباتی بسیار پایین‌تر امکان‌پذیر می‌سازد.

\subsection{مبانی نظری و فرمول‌بندی ریاضی}
فرض کنید ماتریس وزن $\mathbf{W} \in \mathbb{R}^{d_{\text{out}} \times d_{\text{in}}}$ داده شده باشد. ابعاد ورودی و خروجی به عوامل متعادل فاکتورگیری می‌شوند:
\begin{equation}
d_{\text{out}} = m_1 \cdot m_2, \quad d_{\text{in}} = n_1 \cdot n_2
\end{equation}
که در آن ابعاد به گونه‌ای انتخاب می‌شوند که اختلاف $|m_1 - m_2|$ و $|n_1 - n_2|$ کمینه گردد. سپس ماتریس وزن با حاصل‌ضرب کرونکر دو ماتریس فاکتور $\mathbf{A} \in \mathbb{R}^{m_1 \times n_1}$ و $\mathbf{B} \in \mathbb{R}^{m_2 \times n_2}$ بازنمایی می‌شود:
\begin{equation}
\mathbf{W} \approx \mathbf{A} \otimes \mathbf{B} = \begin{bmatrix} 
A_{11}\mathbf{B} & \cdots & A_{1n_1}\mathbf{B} \\ 
\vdots & \ddots & \vdots \\ 
A_{m_11}\mathbf{B} & \cdots & A_{m_1n_1}\mathbf{B} 
\end{bmatrix}
\end{equation}
تعداد پارامترهای وزن از $d_{\text{out}} d_{\text{in}} = m_1 m_2 n_1 n_2$ به $m_1 n_1 + m_2 n_2$ کاهش می‌یابد که کاهش مرتبه‌ای عظیمی را نشان می‌دهد.

برای محاسبه ضرب پیشرو $\mathbf{y} = \mathbf{x} (\mathbf{A} \otimes \mathbf{B})^\top + \mathbf{b}$، از مشخصه دوخطی ضرب کرونکر استفاده می‌شود تا محاسبات بدون تشکیل ماتریس کامل $\mathbf{W}$ انجام گیرد:
\begin{enumerate}
    \item ورودی $\mathbf{x} \in \mathbb{R}^{K \times d_{\text{in}}}$ به ابعاد $(K \cdot n_1, n_2)$ تغییر شکل یافته و با $\mathbf{B}^\top$ منقبض می‌شود:
    \begin{equation}
    \mathbf{X}_B = \text{reshape}(\mathbf{x}, [K \cdot n_1, n_2]) \, \mathbf{B}^\top \in \mathbb{R}^{(K \cdot n_1) \times m_2}
    \end{equation}
    \item نتیجه ترانهاده شده و با $\mathbf{A}^\top$ ضرب می‌گردد:
    \begin{equation}
    \mathbf{X}_A = \text{reshape}(\mathbf{X}_B^\top, [K \cdot m_2, n_1]) \, \mathbf{A}^\top \in \mathbb{R}^{(K \cdot m_2) \times m_1}
    \end{equation}
    \item تانسور حاصل به شکل $(K, m_1 \cdot m_2)$ بازآرایی شده و بایاس به آن افزوده می‌شود.
\end{enumerate}

در لایه‌هایی که تجزیه غیربدیهی کرونکر امکان‌پذیر نباشد (نظیر ابعاد اول)، ماژول به صورت خودکار به تجزیه کم‌رتبه بهینه‌سازی‌شده \lr{SVD} با رتبه انتخابی سوییچ می‌کند.

\subsection{معماری و پیاده‌سازی}
ماژول در فایل \lr{Codes/framework/compression\_engine/kronecker\_compressor.py} پیاده‌سازی شده و از مؤلفه‌های زیر تشکیل یافته است:
\begin{itemize}
    \item \textbf{تابع \lr{factorize\_dim}:} تجزیه بهینه بعد به دو فاکتور متعادل با استفاده از جذر صحیح و جستجوی مقسوم‌علیه مشترک.
    \item \textbf{کلاس \lr{KroneckerLinear}:} ماژول اختصاصی \lr{PyTorch} که جایگزین لایه‌های استاندارد می‌شود. این لایه در صورت تجزیه‌پذیری، پارامترهای فاکتور $\mathbf{A}$ و $\mathbf{B}$ را نگهداری کرده و از تابع \lr{\_forward\_kron} برای ضرب بدون ماتریسی‌سازی صریح بهره می‌گیرد. همچنین برای فاز آموزش تکمیلی، مجهز به شاخه‌های موازی انطباقی \lr{LoRA} است.
    \item \textbf{کلاس \lr{KroneckerCompressor}:} کلاس اصلی مدیریت فشرده‌سازی در سطح مدل که لایه‌های مجاز را شناسایی کرده، خطای فروبنیوس بازسازی را پایش نموده و در صورت فراتر رفتن خطا از حد آستانه، لایه اصلی را بدون تغییر حفظ می‌کند.
\end{itemize}

\subsection{نتایج تجربی}
ارزیابی ماژول \lr{kronecker\_compressor} بر روی مدل‌های زبانی و انکودرهای پایه بر اساس داده‌های تجربی تاییدشده در \lr{thesis\_final/master\_verified\_results.json} انجام شده و در جدول \ref{tab:benchmark_kronecker_compressor} گزارش شده است:

\begin{table}[H]
\centering
\small
\caption{نتایج تجربی و اعتبارسنجی بنچ‌مارک فشرده‌ساز \lr{kronecker\_compressor} روی مدل‌های پایه}
\label{tab:benchmark_kronecker_compressor}
\begin{tabular}{lccccc}
\toprule
\textbf{مدل پایه} \lr{(Base Model)} & \textbf{دقت} \lr{(Accuracy)} & \textbf{شاخص} \lr{(F1-Score)} & \textbf{حافظه} \lr{(Peak VRAM)} & \textbf{توان} \lr{(Throughput)} & \textbf{وضعیت} \\
\midrule
\lr{DistilBERT} & \lr{90.94\%} & \lr{0.9123} & \lr{315.07 MB} & \lr{157,132.82 tok/s} & \lr{PASS} \\
\lr{RoBERTa} & \lr{70.07\%} & \lr{0.7603} & \lr{301.72 MB} & \lr{109,075.39 tok/s} & \lr{PASS} \\
\lr{GPT-2} & \lr{60.32\%} & \lr{0.7038} & \lr{643.30 MB} & \lr{69,983.27 tok/s} & \lr{PASS} \\
\lr{Qwen} & \lr{50.92\%} & \lr{---} & \lr{---} & \lr{---} & \lr{PASS} \\
\bottomrule
\end{tabular}
\end{table}

بر اساس شواهد تجربی ثبت‌شده، این فشرده‌ساز در مدل انکودر سبک \lr{DistilBERT} عملکرد بالا با دقت \lr{90.94\%}، امتیاز \lr{F1} معادل \lr{0.9123} و توان محاسباتی خیره‌کننده \lr{157,132.82 tok/s} به همراه مصرف ناچیز حافظه \lr{315.07 MB} نشان داده است. در مدل‌های بزرگ‌تر، پایداری در سطح قبولی آزمون حفظ شده و در مدل‌های \lr{RoBERTa}، \lr{GPT-2} و \lr{Qwen} به ترتیب دقت‌های \lr{70.07\%}، \lr{60.32\%} و \lr{50.92\%} ثبت شده است.

\subsection{محدودیت‌ها و تحلیل شکست}
با وجود بازدهی و سرعت بالا، روش تجزیه کرونکر دارای محدودیت‌های تئوریک و عملیاتی زیر است:
\begin{itemize}
    \item \textbf{انعطاف‌ناپذیری ابعادی در ابعاد اول:} کیفیت فاکتورگیری کرونکر شدیداً وابسته به تجزیه‌پذیری ابعاد $d_{\text{in}}$ و $d_{\text{out}}$ به فاکتورهای غیربدیهی نزدیک به جذر است. در لایه‌هایی با ابعاد عدد اول یا فاکتورهای نامتعادل، فشردگی کاهش یافته یا لایه ناگزیر به مسیر رزرو \lr{SVD} هدایت می‌شود.
    \item \textbf{افت دقت در دیکودرهای خودبازگشتی طویل:} در مدل‌های زبانی مانند \lr{GPT-2} و \lr{Qwen}، ماهیت ساختار کرونکر قیدهای صلب بلوکی روی درایه‌های وزن ایجاد می‌کند؛ در نتیجه در غیاب تیونینگ کافی، خطای تجمعی در گام‌های توالی طولانی سبب افت کیفیت پیش‌بینی لایه \lr{LM Head} می‌شود.
    \item \textbf{سربار جابجایی تانسورها (\lr{Permutation and Reshape Overheads}):} در پیاده‌سازی گذر پیشرو بدون ماتریسی‌سازی صریح، نیاز به عملیات ترانهاده و تغییر شکل مکرر داده‌ها در حافظه کش وجود دارد که در توالی‌های کوتاه می‌تواند مانع از رسیدن به حداکثر شتاب تئوریک سخت‌افزاری گردد.
\end{itemize}
